\documentclass[11pt,a4paper]{article}
\usepackage[T1]{fontenc}
\usepackage[utf8]{inputenc}
\usepackage{lmodern}
\usepackage[czech]{babel}
\usepackage[margin=2cm]{geometry}
\usepackage{amsmath,amssymb,mathtools}
\usepackage{xcolor,cancel,enumitem}
\usepackage{tikz,pgfplots}
\pgfplotsset{compat=1.18}
\usetikzlibrary{arrows.meta,calc,patterns,shapes.geometric,babel}
\pagestyle{empty}
\setlength{\parindent}{0pt}
\setlength{\parskip}{4pt}
\newcommand{\vysl}[1]{\boxed{\color{red!75!black}#1}}
\newcommand{\pozn}[1]{{\color{red!75!black}#1}}
\newcommand{\dx}{\,\mathrm{d}x}
\newcommand{\dt}{\,\mathrm{d}t}
\newcommand{\subst}[1]{\left|\begin{array}{l}#1\end{array}\right|}
\DeclareMathOperator{\arctg}{arctg}
\DeclareMathOperator{\arccotg}{arccotg}
\DeclareMathOperator{\tg}{tg}
\DeclareMathOperator{\cotg}{cotg}
\usepackage{tcolorbox}
\newtcolorbox{pomocny}{colframe=red!75!black, colback=white, coltitle=white, fonttitle=\bfseries, title={Pomocný výpočet:}, sharp corners}
\newcommand{\dy}{\,\mathrm{d}y}
% znaménkové schéma: \znam{bod/popis/typ,...}{znaménka v intervalech}
% typ: o = vyloučený bod (kroužek), x = obyčejný bod
\newcommand{\znam}[2]{%
\begin{tikzpicture}[baseline=-3pt]
\def\n{0}
\foreach \p/\l/\t [count=\i] in {#1} {\xdef\n{\i}}
\pgfmathsetmacro{\L}{1.6*(\n+1)}
\draw (0,0) -- (\L,0);
\foreach \p/\l/\t [count=\i] in {#1} {
  \pgfmathsetmacro{\xx}{1.6*\i}
  \ifnum\pdfstrcmp{\t}{o}=0 \draw[fill=white] (\xx,0) circle (2.5pt); \draw (\xx-0.09,0) -- (\xx+0.09,0) (\xx,-0.09) -- (\xx,0.09);
  \else \draw (\xx,0.1) -- (\xx,-0.1); \fi
  \node[below=2pt] at (\xx,0) {\small $\l$};
}
\foreach \s [count=\i] in {#2} {
  \pgfmathsetmacro{\xx}{1.6*\i-0.8}
  \node[above] at (\xx,0) {$\s$};
}
\end{tikzpicture}}
\newcommand{\aug}[1]{\left(\begin{array}{rrr|rrr}#1\end{array}\right)}
\newcommand{\krok}[1]{\text{\small\color{orange!80!black}#1}}
% značky pro šipky (více průchodů pdflatex)
\newcommand{\tm}[2]{\tikz[remember picture, baseline=(#1.base)]\node[inner sep=0pt, outer sep=1pt] (#1) {$\displaystyle #2$};}
\newcommand{\tmt}[2]{\tikz[remember picture, baseline=(#1.base)]\node[inner sep=0pt, outer sep=1pt] (#1) {$\textstyle #2$};}
\newcommand{\kruh}[1]{\tikz[baseline=(k.base)]\node[draw, ellipse, inner sep=0pt] (k) {$#1$};}
\newcommand{\mk}[1]{\tikz[remember picture, overlay]\coordinate (#1) at (0,0.25em);}
\tikzset{sip/.style={-{Stealth}, semithick}}
% vodorovná šipka od konce řádku (značka) k hodnotě v bodě
\newcommand{\hsip}[2]{\draw[sip, red!75!black] ([xshift=6pt]#1) -- ([xshift=-5pt]#2.west |- #1);}

\begin{document}

\section*{Maclaurinův polynom 3. stupně\\ \mk{pro}pro $f(x)=\arcsin x$}

\[
T_3(x) = f(0)+\frac{f'(0)\cdot x}{1!}+\frac{f''(0)\cdot x^2}{2!}+\frac{f'''(0)\cdot x^3}{3!}
\]
\hspace*{1.2em}\tm{x0}{x_0=0}
\begin{tikzpicture}[remember picture, overlay]
\draw[sip] ([xshift=-3pt, yshift=1pt]pro) to[out=200, in=180, looseness=1.3] ([xshift=-2pt]x0.west);
\end{tikzpicture}
\begin{align*}
f(0) &= \arcsin0 = 0\\
f'(x) &= \frac{1}{\sqrt{1-x^2}}\mk{a1} & \tm{v1}{f'(0)} &= 1\\
f''(x) &= \left(\frac{1}{\sqrt{1-x^2}}\right)' = \left(\left(1-x^2\right)^{-\frac12}\right)'
= -\frac12\left(1-x^2\right)^{-\frac32}\cdot(-2x) = \frac{x}{\sqrt{\left(1-x^2\right)^3}}\mk{a2} & \tm{v2}{f''(0)} &= 0\\[1.3em]
f'''(x) &= \frac{(x)'\cdot\sqrt{\left(1-x^2\right)^3}-x\cdot\tmt{sq}{\left(\sqrt{\left(1-x^2\right)^3}\right)'}}{\left(\sqrt{\left(1-x^2\right)^3}\right)^2} =\\
&= \frac{1\cdot\sqrt{\left(1-x^2\right)^3}-x\cdot\frac32\left(1-x^2\right)^{\frac12}\cdot(-2x)}{\left(1-x^2\right)^3}\mk{a3} & \tm{v3}{f'''(0)} &= 1
\end{align*}
\begin{tikzpicture}[remember picture, overlay]
\hsip{a1}{v1} \hsip{a2}{v2} \hsip{a3}{v3}
\node[anchor=west] (sqn) at ($(sq.north east)+(1.4,0.42)$) {$\left(1-x^2\right)^{\frac32}$};
\draw[sip] ([yshift=1pt]sq.north) to[out=70, in=180] ([xshift=-2pt]sqn.west);
\end{tikzpicture}
\[
T_3(x) = 0+1\cdot x+0\cdot\frac{x^2}{2!}+\frac{1\cdot x^3}{3!} = \vysl{x+\frac{x^3}{6}}
\]

\section*{Taylorův polynom 3. stupně pro $f(x)=\arctg\frac1x$ v $x_0=1$}

\[
T_3(x) = f(1)+\frac{f'(1)(x-1)}{1!}+\frac{f''(1)(x-1)^2}{2!}+\frac{f'''(1)(x-1)^3}{3!}
\]
\begin{align*}
f(1) &= \arctg1 = \frac\pi4\\
f'(x) &= \frac{1}{1+\left(\frac1x\right)^2}\cdot\left(\frac1x\right)' = \frac{1}{1+\frac{1}{x^2}}\cdot(-1)\cdot x^{-2}
= \frac{1}{\frac{x^2+1}{x^2}}\cdot\frac{-1}{x^2} = \frac{-1}{x^2+1}\mk{b1} & \tm{w1}{f'(1)} &= -\frac12\\
f''(x) &= \left(-\frac{1}{x^2+1}\right)' = \left(-1\cdot\left(x^2+1\right)^{-1}\right)' =\\
&= -1\cdot(-1)\cdot\left(x^2+1\right)^{-2}\cdot2x = \frac{2x}{\left(x^2+1\right)^2}\mk{b2} & \tm{w2}{f''(1)} &= \frac{2}{2^2} = \frac12\\
f'''(x) &= \frac{(2x)'\cdot\left(x^2+1\right)^2-2x\cdot\left(\left(x^2+1\right)^2\right)'}{\left(x^2+1\right)^4} =\\
&= \frac{2\left(x^2+1\right)^2-2x\cdot2\left(x^2+1\right)\cdot2x}{\left(x^2+1\right)^4}\mk{b3}
& \tm{w3}{f'''(1)} &= \frac{2\cdot2^2-2\cdot2\cdot2\cdot2}{2^4}\\
&&&= \frac{8-16}{16} = \frac{-8}{16} = -\frac12
\end{align*}
\begin{tikzpicture}[remember picture, overlay]
\hsip{b1}{w1} \hsip{b2}{w2} \hsip{b3}{w3}
\end{tikzpicture}
\[
\vysl{T_3(x) = \frac\pi4+\left(-\frac12\right)(x-1)+\tm{q1}{\frac14}(x-1)^2-\tm{q2}{\frac1{12}}(x-1)^3}
\]
\begin{tikzpicture}[remember picture, overlay]
\node[anchor=north] (k1) at ($(q1.south)+(0,-1.0)$) {$\dfrac{\frac12}{2!}=\dfrac14$};
\node[anchor=north] (k2) at ($(q2.south)+(0.9,-1.0)$) {$\dfrac{-\frac12}{3!}=\dfrac{-\frac12}{6}=-\dfrac1{12}$};
\draw[sip] (k1.north) -- ([yshift=-5pt]q1.south);
\draw[sip] (k2.north -| q2.south) -- ([yshift=-5pt]q2.south);
\end{tikzpicture}
\vspace{2.2cm}

\newpage
\section*{Určete limity:}

\textbf{a)}
\begin{align*}
\lim_{x\to\frac\pi2}\frac{\sin2x-2\cos x}{\cos^2x}
&= \left[\frac{\sin\bigl(\overbrace{\textstyle 2\cdot\frac\pi2}^{\pi}\bigr)-2\cos\frac\pi2}{\cos^2\frac\pi2} = \frac{0-0}{0} = \frac00\right]
\overset{\text{L.P.}}{=} \lim_{x\to\frac\pi2}\frac{\cos2x\cdot2-2(-\sin x)}{2\cos x\cdot(-\sin x)} =\\
&= \left[\frac{\bigl(\cos\overbrace{\textstyle 2\cdot\frac\pi2}^{\pi}\bigr)\cdot2+2\sin\frac\pi2}{2\cos\frac\pi2\cdot\left(-\sin\frac\pi2\right)}
= \frac{2\cos\pi+2\sin\frac\pi2}{-2\cos\frac\pi2\sin\frac\pi2} = \frac{-2+2}{-2\cdot0\cdot1}\right] = \left[\frac00\right] \overset{\text{L.P.}}{=}\\
&= \lim_{x\to\frac\pi2}\frac{2(-\sin2x)\cdot2+2\cos x}{-2\underbrace{\left(-\sin x\cdot\sin x+\cos x\cdot\cos x\right)}_{(\cos x\sin x)'}} =\\
&= \frac{-2\bigl(\sin\overset{\pi}{\kruh{2\cdot\frac\pi2}}\bigr)\cdot2+2\cos\frac\pi2}{-2\left(-\sin\frac\pi2\cdot\sin\frac\pi2+\cos\frac\pi2\cdot\cos\frac\pi2\right)}
= \frac{-2\cdot0\cdot2+2\cdot0}{-2\cdot(-1\cdot1+0\cdot0)} = \frac02 = \vysl{0}
\end{align*}

\textbf{b)}

\hspace*{3em}$\displaystyle
\lim_{x\to0}\frac{\tmt{lnq}{\ln\frac{\sin x}{x}}}{1-\cos x} = \frac{\ln1}{1-\cos0} = \frac{0}{1-1} = \left[\frac00\right] \ \text{L.P.}$

\medskip
\hspace*{3em}\tm{bx}{\boxed{\lim_{x\to0}\frac{\sin x}{x} = \left[\frac00\right] = \lim_{x\to0}\frac{\cos x}{1} = 1}}
\begin{tikzpicture}[remember picture, overlay]
\node[draw, ellipse, inner sep=0pt, minimum height=0.8cm, minimum width=1.6cm] (lne) at ([yshift=1pt]lnq) {};
\draw[sip] ([yshift=1pt]lne.north west) .. controls +(-1.2,0.6) and +(-1.4,0) .. ([xshift=-3pt]bx.west);
\end{tikzpicture}
\begin{align*}
&\overset{\text{L.P.}}{=} \lim_{x\to0}\frac{\frac{1}{\frac{\sin x}{x}}\cdot\left(\frac{\sin x}{x}\right)'}{0-(-\sin x)}
= \lim_{x\to0}\frac{\frac{1}{\frac{\sin x}{\cancel{x}}}\cdot\frac{\cos x\cdot x-\sin x\cdot1}{\cancel{x^2}}}{\sin x}
= \lim_{x\to0}\frac{\cos x\cdot x-\sin x}{x\cdot\sin^2x} = \left[\frac00\right] \overset{\text{L.P.}}{=}\\
&= \lim_{x\to0}\frac{(-\sin x\cdot x+\cos x)-\cos x}{1\cdot\sin^2x+x\cdot2\sin x\cos x}
= \lim_{x\to0}\frac{-\sin x\cdot x}{\sin^2x+x\cdot2\sin x\cos x} = \lim_{x\to0}\frac{-x}{\sin x+2x\cos x} =\\
&= \left[\frac00\right] \overset{\text{L.P.}}{=} \lim_{x\to0}\frac{-1}{\underbrace{\cos x}_{1}+\underbrace{2\cos x}_{2\cdot1}+\underbrace{2x\cdot(-\sin x)}_{0}} = \vysl{-\frac13}
\end{align*}

\newpage
\section*{Asymptoty:}

\textbf{\textcircled{1} vodorovná pro $f(x)=e^{\frac1x}$}
\begin{align*}
\lim_{x\to\infty}e^{\frac1x} &= e^{\frac1\infty\to0} = e^0 = 1\\
\lim_{x\to-\infty}e^{\frac1x} &= e^{\frac1{-\infty}\to0} = e^0 = 1
\end{align*}
\[ \vysl{y=1 \text{ pro } x\to\pm\infty} \]

\hrule
\medskip
\textbf{\textcircled{2} svislá pro $f(x)=\tm{fl}{\frac{1}{\ln x}}$}
\begin{tikzpicture}[remember picture, overlay]
\node[anchor=west] (n1) at ($(fl.east)+(1.4,0.25)$) {$\ln x\ \dots\ x\in(0,\infty)$};
\node[anchor=north west, align=left] (n2) at ($(fl.east)+(1.6,-0.35)$) {$\ln x\ne0$\\ $\phantom{\ln}\,x\ne1$};
\draw[sip] ([xshift=3pt]fl.east) -- ([xshift=-2pt]n1.west);
\draw[sip] ([xshift=3pt, yshift=-3pt]fl.east) -- ([xshift=-2pt]n2.north west);
\end{tikzpicture}

\medskip
\hspace*{1.5em}$D(f)=(0,1)\cup(1,\infty)$
\[
\lim_{x\to1}\frac{1}{\ln x} = \left\langle\begin{array}{l}
\displaystyle\lim_{x\to1^+}\frac{1}{\ln x} = \frac{1}{0^+} = \infty\\[10pt]
\displaystyle\lim_{x\to1^-}\frac{1}{\ln x} = \frac{1}{0^-} = -\infty
\end{array}\right\}\quad \vysl{x=1}
\]

\hrule
\medskip
\textbf{\textcircled{3} šikmá pro $f(x)=\dfrac{x^2+1}{x}$}
\begin{align*}
a &= \lim_{x\to\infty}\frac{f(x)}{x} = \lim_{x\to\infty}\frac{\frac{x^2+1}{x}}{x} = \lim_{x\to\infty}\frac{x^2+1}{x^2}
= \left[\frac\infty\infty\right] \overset{\text{L.P.}}{=} \lim_{x\to\infty}\frac{2x}{2x} = 1\\
b &= \lim_{x\to\infty}\big(f(x)-ax\big) = \lim_{x\to\infty}\left(\frac{x^2+1}{x}-x\right)
= \lim_{x\to\infty}\frac{x^2+1-x^2}{x} = \lim_{x\to\infty}\frac1x = \frac1\infty = 0
\end{align*}
\[
y=1\cdot x+0 \ \Rightarrow\ \vysl{y=x} \quad \boxed{\begin{array}{l}\text{pro } x\to\infty\\ \text{pro } x\to-\infty\end{array}}
\]
\underline{pozn.} $\displaystyle\lim_{x\to-\infty}$ \dots\ stejný výpočet

\end{document}
